\pagenumbering{Roman}
\begin{titlepage}
  \thispagestyle{firststyle}
  \centering
  \vspace*{22mm}
  {\sffamily\bfseries\color{Navy}\fontsize{29}{34}\selectfont
    Олимпиадная теория чисел\par}
  \vspace{5mm}
  {\sffamily\fontsize{18}{23}\selectfont
    Оценки, показатели, степени вхождения\par}
  \vspace{12mm}

  \vspace{12mm}
  \vfill
  \vspace*{18mm}
\end{titlepage}

\clearpage
\pagenumbering{arabic}
\thispagestyle{plain}
\begingroup
\setlength{\parskip}{0pt}
\tableofcontents
\endgroup

\chapter{Оценочная теория чисел}

Оценочная теория чисел отвечает не только на вопрос
«делится или не делится», но и на вопрос «насколько велико» или
«насколько близко». В олимпиадных листках под этим названием встречаются
два тесно связанных слоя. В первом делимость совмещают с неравенством:
ненулевое кратное не может быть меньше делителя. Во втором принцип
Дирихле даёт \emph{хорошие} рациональные приближения, а
целочисленность и алгебраичность запрещают приближениям быть
\emph{слишком хорошими}.

\section{Основные обозначения и принцип зажатия}

\begin{definitionbox}{Целая и дробная части}
Для $x\in\R$ обозначим
\[
  \lfloor x\rfloor=\max\{n\in\Z:n\le x\},
  \qquad
  \{x\}=x-\lfloor x\rfloor\in[0,1).
\]
Расстояние от $x$ до ближайшего целого обозначается
\[
  \dist{x}=\min_{n\in\Z}|x-n|
  =\min\bigl(\{x\},1-\{x\}\bigr).
\]
В частности, $\dist{x}=0$ тогда и только тогда, когда $x$ целое.
\end{definitionbox}

\begin{definitionbox}{Рациональные и алгебраические числа}
Дробь $p/q$ называется \emph{несократимой}, если $q>0$ и
$(p,q)=1$. Действительное число $\alpha$ называется
\emph{алгебраическим степени $d$}, если $d$ --- наименьшая степень
ненулевого многочлена с целыми коэффициентами, имеющего корень
$\alpha$. Неалгебраические числа называются \emph{трансцендентными}.
\end{definitionbox}

\begin{methodbox}{Целое между $-1$ и $1$}
Если удалось построить целое число $T$ и доказать $|T|<1$, то
$T=0$. Чуть более общая форма: если $T\ne0$ целое, то $|T|\ge1$.
Именно эта простая мысль даёт нижние оценки вида
\[
  \left|\frac ab-\frac pq\right|
  =\frac{|aq-bp|}{bq}\ge\frac1{bq}.
\]
Перед применением надо отдельно убедиться, что $aq-bp\ne0$.
\end{methodbox}

\section{Дискретность делимости и малый остаток}

\begin{theorembox}{Лемма о наименьшем ненулевом кратном}
Если $d,A\in\Z$, $d\ne0$ и $d\mid A$, то либо $A=0$, либо
\[
  |A|\ge|d|.
\]
Следовательно,
\[
  d\mid A,\quad |A|<|d|
  \quad\Longrightarrow\quad A=0.
\]
\end{theorembox}

Доказательство состоит в записи $A=dk$: если $A\ne0$, то
$k\ne0$ и $|k|\ge1$. Строгое неравенство существенно: при
$|A|=|d|$ возможны $A=d$ и $A=-d$.

\begin{methodbox}{Уменьшение делимого}
Из $d\mid A$ следует $d\mid A-td$ для любого $t\in\Z$. Поэтому
вместо большого выражения $A$ надо искать кратное $td$, уничтожающее
его главную часть, и оценивать малый остаток
\[
  R=A-td.
\]
Если $|R|<|d|$, то делимость заставляет $R$ обратиться в нуль.
Для двух выражений
\[
  R=aM+r,\qquad B=bM+s
\]
полезна комбинация $bR-aB=br-as$: большая часть $M$ исчезает.
\end{methodbox}

\begin{theorembox}{Простейшие оценки делителей}
Для натурального $n$ верны следующие факты.
\begin{enumerate}
  \item Если $d\mid n$ и $1\le d<n$, то $d\le n/2$.
  \item Если $n$ составное, то у него есть простой делитель
        $p\le\sqrt n$.
  \item Если $d<d'$ --- положительные делители $n$, то
        \[
          d'-d\ge\frac{dd'}n>\frac{d^2}n.
        \]
\end{enumerate}
\end{theorembox}

В третьем пункте числа $n/d$ и $n/d'$ --- различные целые, поэтому
\[
  1\le\frac nd-\frac n{d'}
  =\frac{n(d'-d)}{dd'}.
\]
Такие оценки особенно полезны после упорядочивания переменных или
делителей.

Полезно различать два масштаба. При фиксированном знаменателе $q$
можно округлить $q\alpha$ до ближайшего целого и получить ошибку
не более $1/(2q)$. Теорема Дирихле утверждает больше: среди
\emph{нескольких} знаменателей найдётся такой, для которого ошибка
имеет порядок $1/q^2$.

\section{Теорема Дирихле о рациональных приближениях}

\begin{theorembox}{Теорема Дирихле}
Для любого $\alpha\in\R$ и любого натурального $N$ существуют
$p\in\Z$ и $q\in\N$ такие, что
\[
  1\le q\le N,
  \qquad
  |q\alpha-p|<\frac1N.
\]
Равносильная запись:
\[
  \left|\alpha-\frac pq\right|<\frac1{Nq}.
\]
\end{theorembox}

\textbf{Доказательство.}
Рассмотрим $N+1$ дробную часть
\[
  \{0\alpha\},\{1\alpha\},\ldots,\{N\alpha\}
\]
и разобьём полуинтервал $[0,1)$ на $N$ полуинтервалов длины $1/N$.
По принципу Дирихле две дробные части, скажем с номерами $i<j$,
попадут в один полуинтервал. Поэтому
\[
  |(j-i)\alpha-p|<\frac1N
\]
для некоторого целого $p$. Остаётся положить $q=j-i$. \hfill$\square$

\begin{theorembox}{Следствие для иррационального числа}
Если $\alpha$ иррационально, то существует бесконечно много
несократимых дробей $p/q$, для которых
\[
  \left|\alpha-\frac pq\right|<\frac1{q^2}.
\]
\end{theorembox}

\textbf{Доказательство.}
Применяем теорему при $N=1,2,3,\ldots$. Поскольку $q\le N$,
\[
  \left|\alpha-\frac pq\right|<\frac1{Nq}\le\frac1{q^2}.
\]
Знаменатели не могут оставаться ограниченными: иначе среди конечного
числа значений $|q\alpha-p|>0$ существовал бы положительный минимум,
противоречащий оценке $|q\alpha-p|<1/N$. После сокращения дробей
получаем бесконечно много различных приближений. \hfill$\square$

\begin{commentbox}{Две теоремы Дирихле}
Здесь используется теорема Дирихле \emph{о приближении действительных
чисел рациональными}. Её не следует путать с другой теоремой
Дирихле --- о бесконечности простых чисел в арифметической прогрессии.
\end{commentbox}

\section{Одновременное приближение и плотность дробных частей}

\begin{theorembox}{Одновременная теорема Дирихле}
Пусть $\alpha_1,\ldots,\alpha_s\in\R$ и $N\in\N$. Тогда существуют
$q,p_1,\ldots,p_s\in\Z$ такие, что
\[
  1\le q\le N^s,
  \qquad
  |q\alpha_i-p_i|<\frac1N
  \quad (i=1,\ldots,s).
\]
\end{theorembox}

\textbf{Доказательство.}
Точки
\[
  (\{k\alpha_1\},\ldots,\{k\alpha_s\}),
  \qquad k=0,1,\ldots,N^s,
\]
лежат в единичном $s$-мерном кубе. Разобьём каждое ребро на $N$
частей; получится $N^s$ малых кубиков. Две из $N^s+1$ точек лежат
в одном кубике. Разность их номеров и даёт $q$, а разности координат
дают требуемые оценки. \hfill$\square$

\begin{theorembox}{Плотность иррационального вращения}
Если $\alpha$ иррационально, то множество дробных частей
\[
  \{\{n\alpha\}:n\in\N\}
\]
плотно в $[0,1]$: каждый непустой интервал внутри $[0,1]$ содержит
дробную часть некоторого положительного кратного $\alpha$.
\end{theorembox}

\textbf{Идея доказательства.}
Для заданного интервала длины $\ell$ теорема Дирихле даёт $q$ с
$0<\dist{q\alpha}<\ell$. Если $\{q\alpha\}=\delta<\ell$, то точки
$0,\delta,2\delta,\ldots$ до первого перехода через $1$ идут с шагом
меньше длины интервала. Если $\{q\alpha\}=1-\delta$, те же точки идут
с таким шагом в обратную сторону. В обоих случаях одна из дробных
частей $\{jq\alpha\}$ попадёт в заданный интервал. Хвост любой такой
последовательности снова плотен, поэтому попаданий на самом деле
бесконечно много.

\section{Нижние оценки: рациональные и алгебраические числа}

\begin{theorembox}{Рациональное число нельзя приблизить слишком хорошо}
Пусть $\alpha=a/b$, где $(a,b)=1$, $b>0$. Если $p/q\ne a/b$ и
$q>0$, то
\[
  \left|\alpha-\frac pq\right|\ge\frac1{bq}.
\]
\end{theorembox}

Действительно, $aq-bp$ --- ненулевое целое число, поэтому
$|aq-bp|\ge1$. Эта оценка имеет порядок $1/q$, тогда как для
иррационального числа теорема Дирихле бесконечно часто даёт порядок
$1/q^2$.

\begin{theorembox}{Теорема Лиувилля}
Пусть $\alpha$ --- действительное алгебраическое число степени
$d>1$. Тогда существует константа $C(\alpha)>0$ такая, что для любых
$p\in\Z$ и $q\in\N$
\[
  \left|\alpha-\frac pq\right|>
  \frac1{C(\alpha)q^d}.
\]
\end{theorembox}

\textbf{Доказательство.}
Пусть $P(x)=a_dx^d+\cdots+a_0$ --- минимальный многочлен числа
$\alpha$. Рациональное число $p/q$ не является его корнем, иначе
$P$ имел бы линейный множитель над $\Q$. Кроме того,
\[
  q^dP\!\left(\frac pq\right)\in\Z\setminus\{0\},
  \qquad
  \left|P\!\left(\frac pq\right)\right|\ge\frac1{q^d}.
\]
Если $|p/q-\alpha|<1$, то на отрезке между $\alpha$ и $p/q$ величина
$|P'(x)|$ ограничена некоторой константой $M$. По теореме Лагранжа
\[
  \frac1{q^d}
  \le\left|P\!\left(\frac pq\right)-P(\alpha)\right|
  \le M\left|\frac pq-\alpha\right|.
\]
Для дробей, удалённых от $\alpha$ хотя бы на $1$, оценка обеспечивается
увеличением общей константы $C(\alpha)$. \hfill$\square$

\begin{commentbox}{Что именно запрещает теорема Лиувилля}
Теорема не утверждает, что алгебраические числа плохо приближаются
рациональными. Теорема Дирихле всё равно даёт ошибку $<1/q^2$.
Лиувилль запрещает для числа степени $d$ бесконечную серию ошибок,
существенно меньших $1/q^d$.
\end{commentbox}

\section{Лемма Зигеля: малая целочисленная зависимость}

\begin{theorembox}{Элементарная лемма Зигеля}
Рассмотрим систему $r$ однородных линейных уравнений с $s$ неизвестными
и целыми коэффициентами $a_{ij}$, где $r<s$ и $|a_{ij}|<A$,
$A\ge1$. Тогда у системы есть ненулевое целочисленное решение
$(x_1,\ldots,x_s)$, для которого
\[
  |x_j|<1+(sA)^{\frac r{s-r}}
  \qquad (j=1,\ldots,s).
\]
\end{theorembox}

\textbf{Доказательство.}
Положим
\[
  T=\left\lfloor(sA)^{\frac r{s-r}}\right\rfloor
\]
и подставим в левые части все $(T+1)^s$ наборов
$(u_1,\ldots,u_s)$, где $0\le u_j\le T$. Для фиксированной строки
значение $\sum_j a_{ij}u_j$ пробегает целые числа в интервале длины
меньше $sAT$, поэтому имеет меньше $sA(T+1)$ возможных значений.
У всех $r$ строк вместе меньше
\[
  (sA)^r(T+1)^r
\]
векторов значений. Но
\[
  (T+1)^{s-r}>(sA)^r,
\]
следовательно, наборов неизвестных больше, чем векторов значений.
Два различных набора дают одинаковые левые части; их разность ---
ненулевое решение, причём $|x_j|\le T<1+(sA)^{r/(s-r)}$.
\hfill$\square$

\begin{commentbox}{Смысл леммы}
При числе неизвестных, большем числа уравнений, ненулевое рациональное
решение ожидаемо. Содержание леммы Зигеля в другом: можно выбрать
\emph{целое} решение и заранее ограничить размеры его координат.
Официальные листки ЦПМ по «оценкам в теории чисел» используют именно
эту форму подсчёта образов и прообразов.
\end{commentbox}

\section{Разбор задач}

\begin{problembox}{Суммы чисел из короткого промежутка}
Для натурального $N$ рассмотрели числа
\[
  N^3,N^3+1,\ldots,N^3+N.
\]
Некоторые $a$ чисел покрасили в красный цвет, затем некоторые $b$
из оставшихся --- в синий; обе группы непусты. Сумма красных чисел
делится на сумму синих. Докажите, что $b\mid a$.
\source{листок кружка в Хамовниках «Оценки в теории чисел»; метод сокращения главной части}
\end{problembox}

\begin{solutionbox}
Обозначим суммы через $R$ и $B$. Поскольку все числа отличаются от
$N^3$ не более чем на $N$,
\[
  R=aN^3+r,\qquad B=bN^3+s,
  \qquad 0\le r\le aN,\quad 0\le s\le bN.
\]
По условию $R=tB$ для некоторого $t\in\N$. Тогда
\[
  bR-aB=(bt-a)B=br-as,
\]
так что $B\mid(br-as)$.

При $N\ge2$ группы не пересекаются и непусты, поэтому
$a+b\le N+1$ и, в частности, $a\le N$. Так как $br,as\ge0$,
\[
  |br-as|\le\max(br,as)\le abN\le bN^2<bN^3\le B.
\]
Лемма о наименьшем ненулевом кратном даёт $br-as=0$. Следовательно,
$(bt-a)B=0$, а поскольку $B>0$, имеем $a=bt$ и $b\mid a$.

При $N=1$ есть всего два числа. Тогда из непустоты и непересечения
групп следует $a=b=1$, и утверждение также верно.

\textbf{Полезный комментарий.}
Оценка $|u-v|\le\max(u,v)$ для неотрицательных $u,v$ здесь сильнее
более привычной, но слишком грубой оценки $|u-v|\le u+v$.
\end{solutionbox}

\begin{problembox}{Одновременное приближение двух чисел}
Докажите, что для любых действительных $\alpha,\beta$ и натурального
$N$ существуют целые $p,r$ и натуральное $q\le N^2$ такие, что
\[
  \left|\alpha-\frac pq\right|<\frac1{Nq},
  \qquad
  \left|\beta-\frac rq\right|<\frac1{Nq}.
\]
\source{официальный листок ЦПМ «Оценки в теории чисел» и материал AoPS о рациональных приближениях}
\end{problembox}

\begin{solutionbox}
Рассмотрим $N^2+1$ точку
\[
  (\{k\alpha\},\{k\beta\}),
  \qquad k=0,1,\ldots,N^2,
\]
в единичном квадрате. Разобьём квадрат на $N^2$ клеток со стороной
$1/N$. Две точки с номерами $i<j$ попадут в одну клетку. Для
$q=j-i$ имеем $1\le q\le N^2$, а разности обеих пар координат по
модулю меньше $1/N$. Поэтому для некоторых целых $p,r$
\[
  |q\alpha-p|<\frac1N,
  \qquad
  |q\beta-r|<\frac1N.
\]
После деления на $q$ получаем требуемое.

\textbf{Полезный комментарий.}
В размерности $s$ знаменатель вырастает до $N^s$. Это не случайный
технический показатель: число малых кубиков равно $N^s$.
\end{solutionbox}

\begin{problembox}{Степень, начинающаяся с цифры $7$}
Пусть $a>1$ --- целое число, не являющееся степенью $10$. Докажите,
что бесконечно много степеней $a^n$ начинаются в десятичной записи
с цифры $7$.
\source{ключевой мотив USAMO 2019, задача 3; плотность дробных частей}
\end{problembox}

\begin{solutionbox}
Положим $\alpha=\log_{10}a$. Число $\alpha$ иррационально. В самом
деле, если $\alpha=u/v$ в несократимом виде, то
\[
  a^v=10^u=2^u5^u.
\]
Из основной теоремы арифметики $v$ делит $u$, следовательно $v=1$ и
$a$ является степенью $10$, вопреки условию.

Запишем $n\alpha=k+\theta$, где $k\in\Z$, $0\le\theta<1$. Тогда
\[
  a^n=10^{n\alpha}=10^k\cdot10^\theta.
\]
Число $a^n$ начинается с цифры $7$ тогда и только тогда, когда
\[
  7\le10^\theta<8,
\]
то есть
\[
  \log_{10}7\le\{n\alpha\}<\log_{10}8.
\]
Этот интервал непуст. Плотность дробных частей $\{n\alpha\}$ даёт
попадание в него, а плотность любого хвоста последовательности даёт
бесконечно много попаданий.

\textbf{Полезный комментарий.}
Условие «первая цифра» естественно переводится в условие на дробную
часть десятичного логарифма. Такой переход часто превращает задачу
о цифрах в задачу о рациональных приближениях.
\end{solutionbox}

\begin{problembox}{Делимость целых частей}
Пусть $\rho>1$ --- действительное число. Известно, что для любых
$m,n\in\N$ из $m\mid n$ следует
\[
  \lfloor m\rho\rfloor\mid\lfloor n\rho\rfloor.
\]
Докажите, что $\rho$ --- целое число.
\source{официальный листок ЦПМ «Оценки в теории чисел», задача 5}
\end{problembox}

\begin{solutionbox}
Предположим, что $\rho$ нецелое. Выберем достаточно большое $m$ так,
чтобы
\[
  A=\lfloor m\rho\rfloor>1
\]
и $m\rho$ не было целым. Такой $m$ существует: для иррационального
$\rho$ подходит любое достаточно большое $m$, а для
$\rho=u/v\notin\Z$ достаточно взять большое $m$, не кратное $v$.

Пусть $\delta=\{m\rho\}\in(0,1)$ и
\[
  k=\left\lceil\frac1\delta\right\rceil.
\]
Тогда $1\le k\delta<2$, поэтому $\lfloor k\delta\rfloor=1$. Так как
$m\mid km$, по условию $A$ должно делить
\[
  \lfloor km\rho\rfloor
  =\lfloor k(A+\delta)\rfloor
  =kA+1.
\]
Но $A\mid kA$ и $A>1$, значит $A\nmid kA+1$. Противоречие. Поэтому
$\rho$ целое.

\textbf{Полезный комментарий.}
Главная оценка здесь --- не асимптотика, а точное попадание
$1\le k\delta<2$, превращающее дробную ошибку ровно в добавку $1$.
\end{solutionbox}

\begin{problembox}{Число Лиувилля}
Докажите, что число
\[
  L=\sum_{j=1}^{\infty}10^{-j!}
   =0{,}110001000000000000000001\ldots
\]
трансцендентно.
\source{классическое следствие теоремы Лиувилля; материал AoPS о рациональных приближениях}
\end{problembox}

\begin{solutionbox}
Оборвём сумму после $n$-го слагаемого:
\[
  L_n=\sum_{j=1}^{n}10^{-j!}=\frac{p_n}{q_n},
  \qquad q_n=10^{n!}.
\]
Хвост оценивается геометрической прогрессией:
\[
  0<L-L_n
  =\sum_{j=n+1}^{\infty}10^{-j!}
  <2\cdot10^{-(n+1)!}
  =\frac{2}{q_n^{\,n+1}}.
\]
Сначала исключим рациональность. Если $L=a/b$, где $b>0$, то
$L_n<L$ и
\[
  \left|L-\frac{p_n}{q_n}\right|
  =\frac{|aq_n-bp_n|}{bq_n}\ge\frac1{bq_n}.
\]
При достаточно большом $n$ это противоречит верхней оценке
$2/q_n^{n+1}$.

Предположим, что $L$ алгебраично степени $d>1$. По теореме Лиувилля
для некоторой константы $C>0$ и всех рациональных $p/q$
\[
  \left|L-\frac pq\right|>\frac1{Cq^d}.
\]
Выберем $n>d$ настолько большим, что
\[
  \frac{2}{q_n^{n+1}}<\frac1{Cq_n^d}.
\]
Для $p_n/q_n$ получаем противоречащие верхнюю и нижнюю оценки.
Следовательно, $L$ не алгебраично, то есть трансцендентно.

\textbf{Полезный комментарий.}
Дробь $p_n/q_n$ может сократиться, но это не мешает: её настоящий
знаменатель не превосходит $q_n$, поэтому нижняя оценка Лиувилля
становится только сильнее.
\end{solutionbox}

\section{Советы}

\begin{methodbox}{Рабочий алгоритм}
\begin{enumerate}
  \item Переведите близость к целому в величину $\dist{q\alpha}$ или
        во дробную часть $\{q\alpha\}$.
  \item Посчитайте объекты и ящики. Для $s$ чисел естественный объект
        --- точка в $s$-мерном кубе.
  \item После получения приближения проверьте масштаб ошибки:
        $1/(Nq)$, $1/q^2$ или $1/q^d$ --- это разные утверждения.
  \item Для нижней оценки найдите ненулевое целое: числитель разности
        дробей или число $q^dP(p/q)$.
  \item В задачах с целой частью пишите $x=\lfloor x\rfloor+\{x\}$;
        это часто отделяет делимую часть от малой поправки.
\end{enumerate}
\end{methodbox}

\section{Частые ошибки}

\begin{warningbox}{Что проверять}
\begin{enumerate}
  \item Из плотности рациональных чисел следует лишь возможность сделать
        ошибку малой; она ничего не говорит о размере знаменателя.
  \item Из теоремы Дирихле нельзя для каждого $q$ заключать
        $|\alpha-p/q|<1/q^2$: хорошим является некоторый знаменатель
        среди $1,\ldots,N$.
  \item После сокращения дроби надо заново понимать, какой знаменатель
        стоит в оценке.
  \item В нижней оценке $|aq-bp|\ge1$ сначала нужно доказать
        $aq-bp\ne0$.
  \item Теорема Лиувилля применима к алгебраическому числу известной
        конечной степени, а не к произвольному иррациональному числу.
  \item Равенство $\lfloor kx\rfloor=k\lfloor x\rfloor$ обычно неверно:
        точная формула содержит $\lfloor k\{x\}\rfloor$.
  \item В лемме Зигеля важны одновременно условия: коэффициенты целые,
        система однородна и неизвестных больше, чем уравнений.
\end{enumerate}
\end{warningbox}

\section{Полезные комментарии}

\begin{commentbox}{Когда оценки становятся равенствами}
Олимпиадная оценка особенно сильна, когда интервал между нижней и
верхней границей содержит единственное целое. Поэтому полезно искать
не просто $X<Y$, а конструкции вида $0\le T<1$, $1\le T<2$ или
$M\le T<M+1$ для целого $T$.
\end{commentbox}

\begin{commentbox}{Граница сложности}
Теоремы Дирихле и Лиувилля имеют элементарные доказательства и доступны
в 10--11 классах. Более сильные результаты о рациональных приближениях
(Гурвиц, Рот) полезны как ориентир, но обычно не требуются в школьной
олимпиаде. Лемма Зигеля в приведённой форме также доказывается одним
подсчётом по принципу Дирихле.
\end{commentbox}

\chapter{Показатели по модулю}

В русской олимпиадной традиции \emph{показатель числа по модулю} ---
это мультипликативный порядок. Его не следует путать ни с показателем
степени, ни со степенью вхождения $\nu_p(n)$.

\section{Определение и существование}

\begin{definitionbox}{Показатель (мультипликативный порядок)}
Пусть $m>1$ и $(a,m)=1$. Показателем числа $a$ по модулю $m$
называется наименьшее натуральное $d$, для которого
\[
  a^d\equiv1\pmod m.
\]
Обозначение: $d=\ord_m(a)$.
\end{definitionbox}

Условие $(a,m)=1$ существенно. Если простой $p$ делит и $a$, и $m$,
то $a^d\equiv0\pmod p$ и не может быть сравнимо с $1$ по модулю $m$.

\begin{theorembox}{Функция Эйлера и теорема Эйлера}
Функция Эйлера $\varphi(m)$ равна числу остатков от $1$ до $m$,
взаимно простых с $m$. Для
$m=p_1^{\alpha_1}\cdots p_s^{\alpha_s}$ имеем
\[
  \varphi(m)=m\prod_{i=1}^{s}\left(1-\frac1{p_i}\right).
\]
При $(a,m)=1$ выполняется
\[
  a^{\varphi(m)}\equiv1\pmod m.
\]
Следовательно, $\ord_m(a)$ существует.
\end{theorembox}

\section{Основная теорема и периодичность}

\begin{theorembox}{Основные свойства показателя}
Пусть $d=\ord_m(a)$. Тогда для неотрицательных $k,r,s$
\begin{align*}
  a^k\equiv1\pmod m
  &\quad\Longleftrightarrow\quad d\mid k,\\
  a^r\equiv a^s\pmod m
  &\quad\Longleftrightarrow\quad r\equiv s\pmod d,
\end{align*}
а также
\[
  d\mid\varphi(m).
\]
\end{theorembox}

\textbf{Доказательство.}
Пусть $a^k\equiv1\pmod m$ и $k=qd+r$, где $0\le r<d$. Тогда
\[
  1\equiv a^k=(a^d)^q a^r\equiv a^r\pmod m.
\]
Минимальность $d$ даёт $r=0$, то есть $d\mid k$. Обратное очевидно.
Во втором утверждении можно сократить сравнение на степень $a$, потому
что $(a,m)=1$, и применить первое утверждение к $|r-s|$. Наконец,
теорема Эйлера и первый пункт дают $d\mid\varphi(m)$. \hfill$\square$

Таким образом, последовательность
\[
  1,a,a^2,a^3,\ldots\pmod m
\]
периодична без предпериода, а длина её наименьшего периода равна
$\ord_m(a)$.

\begin{theorembox}{Показатель степени}
Если $d=\ord_m(a)$ и $k\in\N$, то
\[
  \ord_m(a^k)=\frac d{(d,k)}.
\]
\end{theorembox}

Действительно, $(a^k)^e\equiv1\pmod m$ равносильно $d\mid ke$.
Если $d=gd_1$, $k=gk_1$, $g=(d,k)$, то $(d_1,k_1)=1$ и наименьшее
подходящее $e$ равно $d_1$.

\begin{theorembox}{Разложение составного модуля}
Если $u,v>1$, $(u,v)=1$ и $(a,uv)=1$, то
\[
  \ord_{uv}(a)
  =\lcm\bigl(\ord_u(a),\ord_v(a)\bigr).
\]
Если $1<t\mid m$ и $(a,m)=1$, то
$\ord_t(a)\mid\ord_m(a)$.
\end{theorembox}

Первое равенство следует из китайской теоремы об остатках: сравнение
$a^k\equiv1\pmod{uv}$ равносильно двум сравнениям по модулям $u$ и
$v$. Во втором утверждении любое сравнение по модулю $m$ остаётся
верным по модулю $t$.

\section{Простые делители степенных выражений}

\begin{theorembox}{Главная схема}
Пусть $q$ --- простое число.
\begin{enumerate}
  \item Если $q\mid a^n-1$, то
        \[
          \ord_q(a)\mid n,
          \qquad
          \ord_q(a)\mid q-1.
        \]
  \item Если $q$ нечётно, $q\nmid a$ и $q\mid a^n+1$, то
        \[
          \ord_q(a)\mid2n,
          \qquad
          \ord_q(a)\nmid n.
        \]
\end{enumerate}
\end{theorembox}

Особенно важно не останавливаться на словах «показатель делит $n$».
Если удаётся исключить все собственные делители $n$, то
\[
  \ord_q(a)=n,
  \qquad n\mid q-1,
  \qquad q\equiv1\pmod n.
\]
Для разности $a^n-b^n$ при $q\nmid b$ переходят к обратному остатку:
\[
  (ab^{-1})^n\equiv1\pmod q.
\]

\begin{methodbox}{Как вычислять показатель}
Положите $D=\varphi(m)$. Для каждого простого $\ell\mid D$ проверьте
$a^{D/\ell}\pmod m$: если остаток равен $1$, замените $D$ на $D/\ell$
и повторите. Когда дальнейшее деление невозможно, $D=\ord_m(a)$.
\end{methodbox}

\section{Разбор задач}

\begin{problembox}{Вычисление показателя}
Найдите $\ord_{35}(2)$ и решите в неотрицательных целых числах
сравнение
\[
  2^x\equiv29\pmod{35}.
\]
\source{базовая модель листков ЦПМ и Сириуса «Показатели»}
\end{problembox}

\begin{solutionbox}
Так как $35=5\cdot7$,
\[
  \ord_5(2)=4,
  \qquad
  \ord_7(2)=3.
\]
По формуле для взаимно простых модулей
\[
  \ord_{35}(2)=\lcm(4,3)=12.
\]
Кроме того,
\[
  2^6=64\equiv29\pmod{35}.
\]
Следовательно,
\[
  2^x\equiv2^6\pmod{35}
  \quad\Longleftrightarrow\quad
  x\equiv6\pmod{12}.
\]
Все решения:
\[
  \boxed{x=6+12t\quad(t=0,1,2,\ldots)}.
\]

\textbf{Полезный комментарий.}
После нахождения одного подходящего показателя степени не нужно строить
длинную таблицу остатков: все решения автоматически образуют один
класс по модулю порядка.
\end{solutionbox}

\begin{problembox}{Пятые корни единицы}
Пусть $p>5$ --- простое число и
\[
  p\mid n^4+n^3+n^2+n+1.
\]
Докажите, что $5\mid p-1$.
\source{математическая программа Сириуса, листок «Показатели», задача 3}
\end{problembox}

\begin{solutionbox}
Умножим данное выражение на $n-1$:
\[
  (n-1)(n^4+n^3+n^2+n+1)=n^5-1.
\]
Поэтому $n^5\equiv1\pmod p$. При этом $n\not\equiv1\pmod p$:
иначе исходная сумма была бы сравнима с $5$, что невозможно при
$p>5$. Значит,
\[
  \ord_p(n)=5.
\]
По малой теореме Ферма показатель делит $p-1$, откуда
\[
  \boxed{5\mid p-1}.
\]

\textbf{Полезный комментарий.}
После умножения геометрической суммы возникает равенство пятой степени
с единицей. Исключение случая $n\equiv1$ --- обязательный шаг: без него
показатель мог бы равняться $1$.
\end{solutionbox}

\begin{problembox}{Простой делитель суммы степеней}
Пусть $p,q$ --- простые числа, $q>5$, и
\[
  q\mid2^p+3^p.
\]
Докажите, что $q>p$.
\source{математическая программа Сириуса, листок «Показатели», задача 4}
\end{problembox}

\begin{solutionbox}
По модулю $q$ число $3$ обратимо. Положим
\[
  a\equiv2\cdot3^{-1}\pmod q.
\]
Тогда $a^p\equiv-1\pmod q$, поэтому
\[
  \ord_q(a)\mid2p,
  \qquad
  \ord_q(a)\nmid p.
\]
Так как $p$ простое, возможны только значения $2$ и $2p$. Значение
$2$ невозможно: тогда $a\equiv-1\pmod q$, то есть
$2\equiv-3\pmod q$ и $q\mid5$, вопреки $q>5$. Следовательно,
\[
  \ord_q(a)=2p.
\]
Но показатель делит $q-1$, поэтому $2p\mid q-1$ и
\[
  \boxed{q\ge2p+1>p}.
\]

\textbf{Полезный комментарий.}
Для суммы $2^p+3^p$ деление на $3^p$ превращает два основания в одно.
Остаток $-1$ удваивает естественный кандидат на показатель.
\end{solutionbox}

\begin{problembox}{Квадрат делит $2^n+1$}
Найдите все целые $n>1$, для которых
\[
  n^2\mid2^n+1.
\]
\source{IMO 1990, задача 3}
\end{problembox}

\begin{solutionbox}
Число $n$ нечётно. Пусть $p$ --- его наименьший простой делитель и
$d=\ord_p(2)$. Из $2^n\equiv-1\pmod p$ следует
\[
  d\mid2n,
  \qquad d\nmid n,
  \qquad d\mid p-1.
\]
Поскольку $n$ нечётно, $d=2e$ для некоторого $e\mid n$. Кроме того,
$e<p$. Если бы $e>1$, любой простой делитель $e$ был бы делителем
$n$, меньшим $p$. Это противоречит выбору $p$. Значит, $e=1$ и
$d=2$. Тогда $2\equiv-1\pmod p$, откуда $p=3$.

Итак, $3\mid n$. Обозначим $t=\nu_3(n)$. Для нечётного $n$ лемма
LTE даёт
\[
  \nu_3(2^n+1)=\nu_3(2+1)+\nu_3(n)=1+t.
\]
Из $n^2\mid2^n+1$ получаем $2t\le1+t$, поэтому $t=1$.

Предположим, что у $n$ есть другой простой делитель, и пусть $q>3$
--- наименьший из них. Снова положим $d=\ord_q(2)$. Как выше,
$d=2e$, где $e\mid n$ и $e<q$. Все простые делители $e$ меньше $q$;
ими может быть только $3$, причём $3$ входит в $n$ в первой степени.
Следовательно, $e=1$ или $e=3$.

При $e=1$ получаем $d=2$ и $q=3$, противоречие. При $e=3$ имеем
$d=6$, поэтому $2^3$ имеет порядок $2$ по модулю $q$ и
$2^3\equiv-1\pmod q$. Тогда $q\mid9$, снова противоречие. Других
простых делителей у $n$ нет, значит $n=3$. Проверка:
\[
  3^2=2^3+1=9.
\]
Ответ:
\[
  \boxed{n=3}.
\]

\textbf{Полезный комментарий.}
Показатель ограничивает возможные простые делители $n$, а степень
вхождения исключает лишнюю кратность простого $3$. В сложных задачах
эти два инструмента часто работают именно вместе.
\end{solutionbox}

\section{Советы}

\begin{methodbox}{Рабочий алгоритм}
\begin{enumerate}
  \item Если $q\mid a^n-1$, сразу введите $d=\ord_q(a)$ и выпишите
        $d\mid n$, $d\mid q-1$.
  \item Чтобы доказать $d=n$, исключайте собственные делители $n$.
        При простом $n$ обычно достаточно исключить $d=1$.
  \item Для $a^n+1$ используйте одновременно
        $a^n\equiv-1$ и $a^{2n}\equiv1$.
  \item В выражении $a^n-b^n$ делите сравнение на $b^n$, если
        $(b,q)=1$.
  \item При составном модуле разложите его на взаимно простые части и
        соедините показатели через НОК.
  \item Иногда минимальность показателя доказывается оценкой:
        при $M=a^n-1$ и $1\le k<n$ имеем
        $0<a^k-1<M$, поэтому $M\nmid a^k-1$.
\end{enumerate}
\end{methodbox}

\section{Частые ошибки}

\begin{warningbox}{Что проверять}
\begin{enumerate}
  \item $\ord_m(a)$ определён в обычном смысле только при $(a,m)=1$.
  \item Всегда верно лишь $\ord_m(a)\mid\varphi(m)$; равенство бывает
        не для каждого $a$.
  \item Из $q\mid a^n-1$ следует только $\ord_q(a)\mid n$, а не
        автоматическое равенство $\ord_q(a)=n$.
  \item Сокращать множитель в сравнении можно лишь тогда, когда он
        обратим по данному модулю.
  \item Из $d\mid n$ и $d\mid q-1$ нельзя заключать $n\mid q-1$,
        пока не доказано $d=n$.
  \item Для $a^n+1$ отдельно следите за простым $q=2$, где $-1\equiv1$.
  \item Не смешивайте $\ord_m(a)$ и $\nu_p(N)$: первый объект измеряет
        период степеней, второй --- кратность простого множителя.
\end{enumerate}
\end{warningbox}

\section{Полезные комментарии}

\begin{commentbox}{Показатель как переводчик}
Сравнение со степенями превращается в делимость обычных целых чисел:
\[
  a^r\equiv a^s\pmod m
  \quad\Longleftrightarrow\quad
  r\equiv s\pmod{\ord_m(a)}.
\]
Поэтому показатель особенно полезен, когда неизвестная стоит в
показателе степени, а модуль является простым делителем большого
степенного выражения.
\end{commentbox}

\begin{commentbox}{Первообразный корень}
Если $\ord_m(a)=\varphi(m)$, число $a$ называется первообразным корнем
по модулю $m$. Нельзя без доказательства предполагать, что такой корень
существует для любого модуля или что выбранное $a$ является им. Для
большинства олимпиадных задач достаточно делимости порядка на
$\varphi(m)$, без общей теории первообразных корней.
\end{commentbox}

\chapter{Степени вхождения простых чисел}

\section{Определение и основные свойства}

\begin{definitionbox}{Степень вхождения}
Пусть $p$ --- простое число, а $a\ne0$ --- целое. Степенью вхождения
$p$ в $a$, или $p$-адической оценкой числа $a$, называется
\[
  \nu_p(a)=\max\{\alpha\ge0:p^\alpha\mid a\}.
\]
Эквивалентная запись:
\[
  \nu_p(a)=\alpha
  \quad\Longleftrightarrow\quad
  p^\alpha\mid a,
  \quad p^{\alpha+1}\nmid a.
\]
Последнее часто обозначают $p^\alpha\parallel a$. Удобно положить
$\nu_p(0)=+\infty$.
\end{definitionbox}

Например,
\[
  360=2^3\cdot3^2\cdot5,
  \qquad
  \nu_2(360)=3,
  \quad\nu_3(360)=2,
  \quad\nu_5(360)=1.
\]
Знак не влияет на оценку. Для ненулевого рационального числа полагают
\[
  \nu_p\!\left(\frac ab\right)=\nu_p(a)-\nu_p(b);
\]
такая оценка может быть отрицательной.

\begin{definitionbox}{Примарная компонента}
Число
\[
  p^{\nu_p(n)}
\]
называют $p$-примарной компонентой натурального $n$. Степень
$\nu_p(n)$ сообщает кратность, а примарная компонента хранит целиком
наибольшую степень $p$, делящую $n$.
\end{definitionbox}

\begin{theorembox}{Арифметика оценок}
Для ненулевых целых или рациональных $a,b$ и натурального $m$:
\begin{align*}
  \nu_p(ab)&=\nu_p(a)+\nu_p(b),\\
  \nu_p(a^m)&=m\nu_p(a),\\
  \nu_p(a+b)&\ge
  \min\bigl(\nu_p(a),\nu_p(b)\bigr).
\end{align*}
Если $\nu_p(a)\ne\nu_p(b)$, то в последней формуле достигается
равенство. То же верно для разности.
\end{theorembox}

Для доказательства последнего свойства пусть, например,
$\nu_p(a)=r<\nu_p(b)$. Тогда
\[
  a=p^ra_0,
  \qquad b=p^rp^sb_0,
  \qquad s>0,\quad p\nmid a_0.
\]
В скобке $a_0+p^sb_0$ после вынесения $p^r$ первый член не делится
на $p$, а второй делится. Поэтому вся скобка не делится на $p$.

\begin{commentbox}{Единственный минимум}
Если среди нескольких слагаемых ровно одно имеет наименьшую
$p$-адическую оценку, то оценка всей суммы равна этому минимуму.
Увеличение оценки возможно только при «сокращении» нескольких членов
одинаковой минимальной оценки.
\end{commentbox}

\begin{theorembox}{Делимость, НОД, НОК и точные степени}
Для натуральных $a,b$:
\begin{align*}
  a\mid b
  &\quad\Longleftrightarrow\quad
  \nu_p(a)\le\nu_p(b)\ \text{для каждого простого }p,\\
  \nu_p((a,b))&=\min(\nu_p(a),\nu_p(b)),\\
  \nu_p(\lcm(a,b))&=\max(\nu_p(a),\nu_p(b)).
\end{align*}
Натуральное $N$ является точной $m$-й степенью тогда и только тогда,
когда $m\mid\nu_p(N)$ для каждого простого $p$.
\end{theorembox}

Если $M=p_1^{\alpha_1}\cdots p_s^{\alpha_s}$, то наибольшее $r$,
для которого $M^r\mid N$, равно
\[
  \min_{1\le i\le s}
  \left\lfloor\frac{\nu_{p_i}(N)}{\alpha_i}\right\rfloor.
\]
В частности, число нулей в конце десятичной записи $N$ равно
$\min(\nu_2(N),\nu_5(N))$.

\section{Лемма о подъёме показателя (LTE)}

\begin{theorembox}{LTE для нечётного простого}
Пусть $p$ --- нечётное простое,
\[
  p\mid x-y,
  \qquad p\nmid xy,
  \qquad n\in\N.
\]
Тогда
\[
  \boxed{\nu_p(x^n-y^n)=\nu_p(x-y)+\nu_p(n)}.
\]
Если $p\mid x+y$, $p\nmid xy$ и $n$ нечётно, то
\[
  \boxed{\nu_p(x^n+y^n)=\nu_p(x+y)+\nu_p(n)}.
\]
\end{theorembox}

\textbf{Идея доказательства.}
Если $p\nmid m$, то
\[
  \frac{x^m-y^m}{x-y}
  =x^{m-1}+x^{m-2}y+\cdots+y^{m-1}
  \equiv my^{m-1}\not\equiv0\pmod p.
\]
Значит, возведение показателя в множитель, не делящийся на $p$, не
меняет оценку. Далее, если $X=Y+p^tu$ и $p\nmid uY$, раскрытие бинома
даёт
\[
  \nu_p(X^p-Y^p)=t+1=\nu_p(X-Y)+1.
\]
Представив $n=mp^s$, где $p\nmid m$, применяем этот шаг $s$ раз.
Формула для суммы получается заменой $y$ на $-y$; нечётность $n$
сохраняет знак.

\begin{theorembox}{LTE для двойки}
Пусть $x,y$ нечётны.
\begin{enumerate}
  \item Если $n$ нечётно, то
        \[
          \nu_2(x^n-y^n)=\nu_2(x-y),
          \qquad
          \nu_2(x^n+y^n)=\nu_2(x+y).
        \]
  \item Если $n$ чётно, то
        \[
          \boxed{\nu_2(x^n-y^n)
          =\nu_2(x-y)+\nu_2(x+y)+\nu_2(n)-1}.
        \]
  \item Если $n$ чётно, то $\nu_2(x^n+y^n)=1$.
\end{enumerate}
\end{theorembox}

Для второго пункта записывают $n=2^sm$, где $m$ нечётно, и раскладывают
\[
  x^n-y^n=(x^m-y^m)(x^m+y^m)
  \prod_{j=1}^{s-1}(x^{2^jm}+y^{2^jm}).
\]
Начиная с третьего множителя каждая сумма сравнима с $2$ по модулю
$8$ и имеет оценку $1$.

\begin{warningbox}{Проверка условий LTE}
Перед применением леммы явно ответьте на пять вопросов:
\begin{enumerate}
  \item выбранное $p$ действительно простое?
  \item $p\nmid xy$?
  \item делится ли на $p$ именно $x-y$ или $x+y$?
  \item нечётен ли $n$ в формуле для суммы?
  \item не является ли $p=2$, требующим отдельной формулы?
\end{enumerate}
\end{warningbox}

\section{Формула Лежандра}

\begin{theorembox}{Степень простого в факториале}
Для простого $p$
\[
  \boxed{
  \nu_p(n!)=
  \left\lfloor\frac np\right\rfloor+
  \left\lfloor\frac n{p^2}\right\rfloor+
  \left\lfloor\frac n{p^3}\right\rfloor+\cdots}.
\]
Сумма конечна: после $p^j>n$ все члены равны нулю.
\end{theorembox}

Среди $1,\ldots,n$ имеется $\lfloor n/p\rfloor$ кратных $p$.
Кратные $p^2$ содержат ещё по одному множителю $p$, кратные $p^3$ ---
ещё по одному, и так далее.

Пусть $s_p(n)$ --- сумма цифр $n$ в системе счисления по основанию
$p$. Если $n=a_0+a_1p+\cdots+a_rp^r$, то суммирование геометрических
прогрессий даёт цифровую форму:
\begin{theorembox}{Цифровая формула Лежандра}
\[
  \boxed{\nu_p(n!)=\frac{n-s_p(n)}{p-1}}.
\]
\end{theorembox}

\section{Биномиальные коэффициенты и переносы}

Для целых $n\ge0$ и $0\le k\le n$ из формулы Лежандра следует
\[
  \nu_p\binom nk
  =\frac{s_p(k)+s_p(n-k)-s_p(n)}{p-1}.
\]

\begin{theorembox}{Теорема Куммера}
Для целых $n\ge0$ и $0\le k\le n$ число
$\nu_p\binom nk$ равно количеству переносов при сложении
$k$ и $n-k$ в системе счисления с основанием $p$.
\end{theorembox}

Каждый перенос уменьшает сумму цифр результата по сравнению с суммой
цифр слагаемых ровно на $p-1$. Поэтому цифровая формула выше буквально
считает переносы. В частности,
\[
  p\nmid\binom nk
\]
тогда и только тогда, когда сложение $k+(n-k)$ проходит без переносов;
равносильно, при вычитании $k$ из $n$ нет заимствований.

Например,
\[
  \nu_2\binom{2n}{n}
  =2s_2(n)-s_2(2n)=s_2(n),
\]
поскольку умножение на $2$ сдвигает двоичную запись. Поэтому
$\binom{2n}{n}$ делится на $2$, но не на $4$, ровно когда $n$ ---
степень двойки.

\section{Разбор задач}

\begin{problembox}{Наибольшая степень $12$ в факториале}
Найдите наибольшее $r$, для которого
\[
  12^r\mid1000!.
\]
\source{базовая модель применения формулы Лежандра к составному основанию}
\end{problembox}

\begin{solutionbox}
Разложим $12=2^2\cdot3$. По формуле Лежандра
\begin{align*}
  \nu_2(1000!)
  &=500+250+125+62+31+15+7+3+1=994,\\
  \nu_3(1000!)
  &=333+111+37+12+4+1=498.
\end{align*}
Одна степень $12$ требует две двойки и одну тройку, поэтому
\[
  r=\min\left(
  \left\lfloor\frac{994}{2}\right\rfloor,
  498\right)=\boxed{497}.
\]

\textbf{Полезный комментарий.}
Для составного основания нельзя использовать символ $\nu_{12}$ как
обычную простую оценку и складывать показатели. Надо разложить основание
и найти ограничивающую примарную компоненту.
\end{solutionbox}

\begin{problembox}{Локальные степени дают глобальную степень}
Пусть $b,n>1$ --- натуральные числа. Для каждого натурального $k>1$
существует целое $a_k$ такое, что
\[
  k\mid b-a_k^n.
\]
Докажите, что $b$ является точной $n$-й степенью натурального числа.
\source{П. В. Бибиков, листок ЦПМ «Степени вхождения и примарные компоненты», задача 1}
\end{problembox}

\begin{solutionbox}
Пусть $p$ --- любой простой делитель $b$ и $e=\nu_p(b)$. В условии
выберем
\[
  k=p^{e+1}
\]
и обозначим соответствующее число $a_k$ через $a$. Тогда
\[
  \nu_p(b-a^n)\ge e+1.
\]
Если бы $\nu_p(b)\ne\nu_p(a^n)$, свойство оценки разности дало бы
\[
  \nu_p(b-a^n)
  =\min\bigl(e,n\nu_p(a)\bigr)\le e,
\]
противоречие. Следовательно,
\[
  e=n\nu_p(a),
\]
то есть $n\mid e$. Так верно для каждого простого делителя $b$.
Все показатели в разложении $b$ кратны $n$, поэтому
\[
  b=\left(\prod_{p\mid b}p^{\nu_p(b)/n}\right)^n.
\]

\textbf{Полезный комментарий.}
Выбор модуля на одну степень выше, $p^{e+1}$, заставляет две величины
иметь одинаковую оценку. Модуль $p^e$ был бы слишком слаб: он не
различал бы точную кратность.
\end{solutionbox}

\begin{problembox}{Три простые степени одновременно}
Пусть $n$ --- наименьшее натуральное число, для которого
\[
  3^3\cdot5^5\cdot7^7\mid149^n-2^n.
\]
Найдите количество положительных делителей $n$.
\source{AIME I 2020, задача 12; разборы AoPS}
\end{problembox}

\begin{solutionbox}
Так как $149-2=147=3\cdot7^2$, LTE даёт
\begin{align*}
  \nu_3(149^n-2^n)&=1+\nu_3(n),\\
  \nu_7(149^n-2^n)&=2+\nu_7(n).
\end{align*}
Поэтому необходимы и достаточны условия
\[
  3^2\mid n,
  \qquad7^5\mid n.
\]

Для простого $5$ нельзя сразу применить LTE: $5\nmid147$. Сначала
из сравнения $149^n\equiv2^n\pmod5$ получаем
\[
  2^n\equiv1\pmod5,
\]
поэтому $4\mid n$. Положим $n=4c$. Тогда
\[
  149^n-2^n=(149^4)^c-16^c.
\]
Так как $149\equiv-1\pmod{25}$,
\[
  149^4-16\equiv1-16\equiv-15\pmod{25},
\]
следовательно $\nu_5(149^4-16)=1$. Теперь LTE даёт
\[
  \nu_5(149^n-2^n)=1+\nu_5(c).
\]
Для делимости на $5^5$ необходимо и достаточно $5^4\mid c$.

Минимальное число равно
\[
  n=\lcm(3^2,7^5,4\cdot5^4)
   =2^2\cdot3^2\cdot5^4\cdot7^5.
\]
Количество его положительных делителей:
\[
  (2+1)(2+1)(4+1)(5+1)=\boxed{270}.
\]

\textbf{Полезный комментарий.}
Если $p\nmid x-y$, сначала найдите показатель $d$, для которого
$x^d\equiv y^d\pmod p$, и только затем применяйте LTE к новым
основаниям $x^d,y^d$.
\end{solutionbox}

\begin{problembox}{НОД внутренних биномиальных коэффициентов}
Для $n>1$ найдите
\[
  G_n=\gcd\left\{\binom n1,\binom n2,\ldots,\binom n{n-1}\right\}.
\]
\source{классическая олимпиадная задача о теореме Куммера; варианты в материалах AoPS}
\end{problembox}

\begin{solutionbox}
Докажем формулу
\[
  \boxed{
  G_n=
  \begin{cases}
    p,& n=p^a\text{ для некоторого простого }p,\\
    1,& n\text{ не является степенью простого}.
  \end{cases}}
\]

Пусть $n=p^a$. В $p$-ичной записи это единица и $a$ нулей. Для
$0<k<n$ сложение $k+(n-k)$ обязательно содержит перенос, поэтому по
Куммеру $p\mid\binom nk$. Значит, $p\mid G_n$. С другой стороны,
$G_n$ делит $\binom n1=n=p^a$, так что других простых делителей нет.
Чтобы исключить $p^2$, возьмём $k=p^{a-1}$. При сложении
$k+(n-k)$ возникает ровно один перенос, поэтому
\[
  \nu_p\binom{p^a}{p^{a-1}}=1.
\]
Следовательно, $G_n=p$.

Теперь пусть $n$ не является степенью простого. Любой простой делитель
$G_n$ должен делить $\binom n1=n$; пусть это $p$. В $p$-ичной записи
$n$ не является единственной единицей с последующими нулями. Выберем
разряд $j$ с ненулевой цифрой так, чтобы $k=p^j<n$. Все цифры $k$ не
превосходят соответствующих цифр $n$, поэтому при вычитании $k$ из
$n$ нет заимствований, а при сложении $k+(n-k)$ нет переносов.
По Куммеру
\[
  p\nmid\binom nk.
\]
Значит, ни один простой делитель $n$ не делит все внутренние
биномиальные коэффициенты, и $G_n=1$.

\textbf{Полезный комментарий.}
НОД удобно исследовать «по одному простому»: для каждого кандидата
надо найти хотя бы один коэффициент с нулевой оценкой либо доказать
положительную оценку у всех коэффициентов.
\end{solutionbox}

\section{Советы}

\begin{methodbox}{Рабочий алгоритм}
\begin{enumerate}
  \item Переведите $p^a\mid N$ в неравенство $\nu_p(N)\ge a$.
  \item Разложите произведения и дроби: оценки множителей складываются,
        оценки знаменателей вычитаются.
  \item Для суммы сначала сравните оценки слагаемых. Если минимум
        единственный, ответ уже известен.
  \item Для $x^n\pm y^n$ проверьте LTE до раскрытия скобок.
  \item Если нужный простой не делит $x-y$, попробуйте сначала найти
        период по модулю $p$ и сгруппировать показатель.
  \item Для факториала начинайте с Лежандра; для биномиального
        коэффициента сравните этот путь с подсчётом переносов.
  \item После получения необходимых ограничений на неизвестное число
        обязательно проверьте достаточность и минимальность.
\end{enumerate}
\end{methodbox}

\section{Частые ошибки}

\begin{warningbox}{Что проверять}
\begin{enumerate}
  \item $\nu_p(n)$ --- показатель $\alpha$, а не сама степень $p^\alpha$.
  \item Формула
        $\nu_p(a+b)=\min(\nu_p(a),\nu_p(b))$ безусловно неверна;
        равенство гарантировано при разных оценках.
  \item Формулу LTE для нечётного простого нельзя применять к $p=2$.
  \item В LTE нельзя забывать условия $p\nmid xy$ и нужную чётность
        показателя для суммы.
  \item Для составного основания степени простых надо нормировать:
        для $12=2^2\cdot3$ сравниваются $\lfloor\nu_2/2\rfloor$ и
        $\nu_3$.
  \item В формуле Лежандра нельзя останавливаться после
        $\lfloor n/p\rfloor$: кратные $p^2,p^3,\ldots$ дают
        дополнительные множители.
  \item Проверка делимости по одному простому не доказывает делимость
        по всему составному модулю.
\end{enumerate}
\end{warningbox}

\section{Полезные комментарии}

\begin{commentbox}{Локализация по простым и связь с цифрами}
Оценки превращают умножение в сложение, деление --- в вычитание, а
делимость --- в систему неравенств. Поэтому одна глобальная задача
распадается на независимые локальные задачи для простых $p$; ответы
затем собираются основной теоремой арифметики.

Формулы Лежандра и Куммера связывают делимость с записью чисел в
системе счисления по основанию $p$. Высокая степень $p$ в биномиальном
коэффициенте --- это буквально большое число переносов. Такой перевод
часто заменяет громоздкий подсчёт факториалов коротким наблюдением о
цифрах.
\end{commentbox}
